% !TEX root = ../../main.tex

\section{Contexto}

La programación probabilística es un paradigma que permite describir modelos con incertidumbre mediante programas que incorporan operaciones para extraer valores desde distribuciones de probabilidad y, en muchos lenguajes, condicionar esos valores mediante observaciones. A diferencia de un programa determinista, cuyo resultado queda determinado por su entrada, un programa probabilístico especifica implícitamente una distribución sobre sus posibles resultados. La inferencia probabilística consiste en obtener una representación de esa distribución o calcular propiedades de ella, como probabilidades, valores esperados o muestras \cite{GordonHNR2014}.

Como lenguaje de modelado, la programación probabilística se ha aplicado en dominios como visión computacional, teoría de códigos, protocolos criptográficos, biología y análisis de confiabilidad \cite{GordonHNR2014}. Su relación con el aprendizaje automático y la inteligencia artificial es especialmente estrecha: los modelos probabilísticos permiten representar y manipular incertidumbre sobre datos, modelos y predicciones, mientras que los sistemas de programación probabilística permiten expresar esos modelos como procesos computacionales sobre los cuales aplicar inferencia bayesiana \cite{Ghahramani2015,ScibiorGG2015}.

El resultado observado al evaluar un programa probabilístico depende de la estrategia de interpretación utilizada. Bajo una interpretación por muestreo, cada ejecución concreta produce una muestra de la distribución inducida; por ejemplo, una ejecución del lanzamiento de una moneda retorna ``cara'' o ``sello''. Una estrategia de inferencia exacta, en cambio, puede construir una representación explícita de la distribución completa. En ambos casos, el programa describe semánticamente la distribución sobre todos los resultados posibles, no solo una muestra particular \cite{GordonHNR2014,ErwigKollmansberger2006}.

En lenguajes funcionales puros como \textsf{Haskell}, es posible encapsular este tipo de cómputos en tipos de datos y componerlos mediante mónadas de probabilidad. Para distribuciones discretas y finitas, una representación monádica puede construir explícitamente los resultados posibles junto con sus probabilidades, lo que permite realizar inferencia exacta sobre la distribución inducida por el programa \cite{ErwigKollmansberger2006}. El mismo enfoque monádico también puede servir como interfaz para métodos de inferencia más generales, sin restringirse a la enumeración exhaustiva \cite{ScibiorGG2015}.

Estas representaciones permiten realizar consultas y operaciones sobre la distribución de salida de un programa. Por ejemplo, la distribución conjunta asociada al lanzamiento de una moneda seguido del lanzamiento de un dado puede expresarse mediante la \textit{do-notation} de \textsf{Haskell} \cite{ErwigKollmansberger2006}:

\begin{lstlisting}[style=haskell]
    flipCoinRollDice = do 
        c <- uniform [Head, Tail]
        d <- uniform [1 .. 6]
        return (c, d)
\end{lstlisting}

En cada resultado posible del código anterior, la variable \verb|c| contiene un resultado del lanzamiento de la moneda y la variable \verb|d| contiene una cara del dado. Como la segunda elección no depende del resultado de la primera, \verb|return (c, d)| construye la distribución conjunta de ambos experimentos. Su soporte contiene los doce pares obtenidos al combinar los dos resultados de la moneda con las seis caras del dado y, bajo las elecciones uniformes del ejemplo, cada par tiene probabilidad $1/12$ \cite{ErwigKollmansberger2006}.

Tal como se aprecia en el ejemplo, la \textit{do-notation} es especialmente conveniente para programas probabilísticos: un statement de enlace puede introducir una variable aleatoria o un resultado intermedio, y los statements posteriores pueden depender de los valores ya definidos. En el desazucarado monádico convencional de \textsf{Haskell}, un statement de la forma \verb|p <- e| se traduce mediante el operador \textit{bind} \verb|(>>=)|; los statements sin enlace emplean \verb|(>>)| y las declaraciones \texttt{let} poseen una regla propia \cite[sección~3.14]{Haskell2010}. El operador \textit{bind} permite que la forma de un cómputo posterior dependa del valor producido por uno anterior, lo que resulta fundamental para modelar efectos como estado, entrada/salida, fallas, no determinismo o distribuciones probabilísticas \cite{Wadler1995,ErwigKollmansberger2006}.

La secuencia monádica es correcta para el caso general, pero puede ocultar independencia entre subcómputos. Esta tensión fue abordada en \textsf{GHC} mediante la extensión \texttt{ApplicativeDo}, propuesta por Marlow et al. \cite{MarlowPJKM2016}. La extensión analiza las dependencias de un bloque escrito con \textit{do-notation} y, cuando es posible, construye un plan de ejecución mixto que utiliza operadores aplicativos como \verb|(<$>)| y \verb|(<*>)| junto con las operaciones monádicas que todavía sean necesarias \cite{MarlowPJKM2016,GHCApplicativeDo}. De este modo, el programador conserva la sintaxis habitual y el plan resultante puede exponer independencia que ciertas instancias aplicativas aprovechan mediante paralelismo o agrupación de operaciones; la transformación, por sí sola, no garantiza una ejecución multihilo.

Sin embargo, el algoritmo estándar de \texttt{ApplicativeDo} está diseñado para mónadas generales, sin asumir que sus efectos sean conmutativos. Por esta razón, para obtener el nuevo plan de ejecución agrupa statements compatibles pero no cambia su orden relativo: una permutación de estos podría alterar efectos observables, mientras que el intercambio de cómputos independientes solo es válido bajo una hipótesis adicional de conmutatividad \cite[pp.~94--95]{MarlowPJKM2016}. La subsección siguiente muestra, mediante un programa probabilístico concreto, cómo esta restricción puede impedir que el algoritmo de \texttt{ApplicativeDo} genere un plan de ejecución de menor costo.
